\clearpage
\section{Runtime Instruction Examples}

This section shows how compilers and runtime support code can use repeated scalar execution, architectural discovery,
and state-management instructions. The examples and candidate shapes are guidance, not ABI-mandated entry points or
implementation-performance guarantees.

\subsection{Repeat and Streaming Candidates}
Regular memory traversal forms such as MOV, CMP, EXTZ, EXTS, cache maintenance, and selected atomic-free
transforms are natural repeated-execution patterns when combined with a defined repeat instruction or REPG. An implementation may also stream
fixed-size instruction groups when REPG describes a bounded group that has no unresolved control transfer into or
out of the group body.

Copy and fill loops commonly use REPG with MOV over one or two Rn-stepped streams. Scan and compare loops use TEST or
CMP when the grouped body defines termination. Bulk extension combines EXTS or EXTZ with source and destination
auto-update, while cache-maintenance loops traverse regular cache-line sequences. A fixed instruction group is also a
candidate when its byte count ends on an instruction boundary and its body is straight-line. These shapes permit
recognition; they do not guarantee a particular throughput or internal width.

\subsection{Compiler Use}
Compilers may emit repeated scalar forms for compact code size and predictable restart state. They do not need to
know whether a specific processor executes the loop with scalar issue or a streaming engine. Generated code must
still preserve aliasing, ordering, and exception behavior exactly as required by the scalar instruction sequence.

\subsection{CPUID-Defined Program Context Save}
SAVE and RESTORE use a CPUID-described memory image. A runtime that wants to save program-visible register and
extension state first queries the required save-area size and component descriptors, then supplies a
4 KiB-aligned buffer to SAVE.

The continuation state that is not part of the SAVE area, such as a scheduler or setjmp return point, is managed
by the runtime layer that calls this helper.

The two target hooks below are ordinary runtime entry points implemented with the architectural \texttt{CPUID Rn}
and \texttt{SAVE \textless{}ea\textgreater{}} instructions.

\par\begingroup\footnotesize
\begin{verbatim}
#include <stddef.h>
#include <stdint.h>

#define BR_CPUID_CLASS_IMPL       0x00000002u
#define BR_CPUID_SAVE_AREA_LAYOUT 0x0004u
#define BR_SAVE_ALIGNMENT         4096u

struct br_save_area_info {
    uint64_t size;
    uint16_t fixed_size;
    uint16_t component_count;
    uint16_t bitmap_words;
    uint8_t format;
};

struct br_save_component {
    uint16_t id;
    uint8_t bitmap_bit;
    uint32_t offset;
    uint32_t max_size;
    uint16_t alignment;
    uint8_t init_policy;
};

struct br_program_context {
    void *save_area;
    uint64_t save_area_size;
    uint16_t fixed_size;
    uint16_t component_count;
    uint16_t bitmap_words;
    uint8_t format;
};

uint64_t bedrock_cpuid(uint64_t selector);
void bedrock_save(void *save_area);

static uint64_t br_cpuid_selector(uint32_t cls,
                                  uint16_t leaf,
                                  uint16_t index)
{
    return ((uint64_t)cls << 32) |
           ((uint64_t)leaf << 16) |
           (uint64_t)index;
}

static struct br_save_area_info br_query_save_area(void)
{
    uint64_t r0 = bedrock_cpuid(
        br_cpuid_selector(BR_CPUID_CLASS_IMPL,
                          BR_CPUID_SAVE_AREA_LAYOUT, 1));
    uint64_t r1 = bedrock_cpuid(
        br_cpuid_selector(BR_CPUID_CLASS_IMPL,
                          BR_CPUID_SAVE_AREA_LAYOUT, 2));
    struct br_save_area_info info;

    info.size = r0;
    info.fixed_size = (uint16_t)(r1 & 0xffffu);
    info.component_count = (uint16_t)((r1 >> 16) & 0xffffu);
    info.bitmap_words = (uint16_t)((r1 >> 32) & 0xffffu);
    info.format = (uint8_t)((r1 >> 48) & 0x0fu);
    return info;
}

static void br_query_save_components(struct br_save_component *out,
                                     uint16_t count)
{
    for (uint16_t i = 0; i != count; ++i) {
        uint16_t descriptor_a = (uint16_t)(3u + 2u * i);
        uint16_t descriptor_b = (uint16_t)(descriptor_a + 1u);
        uint64_t a = bedrock_cpuid(
            br_cpuid_selector(BR_CPUID_CLASS_IMPL,
                              BR_CPUID_SAVE_AREA_LAYOUT,
                              descriptor_a));
        uint64_t b = bedrock_cpuid(
            br_cpuid_selector(BR_CPUID_CLASS_IMPL,
                              BR_CPUID_SAVE_AREA_LAYOUT,
                              descriptor_b));
        out[i].id = (uint16_t)(a & 0xffffu);
        out[i].bitmap_bit = (uint8_t)((a >> 16) & 0x3fu);
        out[i].offset = (uint32_t)(a >> 32);
        out[i].max_size = (uint32_t)b;
        out[i].alignment = (uint16_t)((b >> 32) & 0xffffu);
        out[i].init_policy = (uint8_t)((b >> 48) & 0xffu);
    }
}

static uintptr_t br_align_up(uintptr_t value, uintptr_t alignment)
{
    return (value + alignment - 1u) & ~(alignment - 1u);
}

int br_save_program_context(struct br_program_context *ctx,
                            struct br_save_component *components,
                            uint16_t component_capacity,
                            void *storage,
                            size_t storage_size)
{
    struct br_save_area_info info = br_query_save_area();
    uintptr_t raw = (uintptr_t)storage;
    uintptr_t base = br_align_up(raw, BR_SAVE_ALIGNMENT);
    size_t leading = (size_t)(base - raw);

    if (leading > storage_size)
        return -1;
    if (info.size > (uint64_t)(storage_size - leading))
        return -1;
    if (info.component_count > component_capacity)
        return -1;

    br_query_save_components(components, info.component_count);

    ctx->save_area = (void *)base;
    ctx->save_area_size = info.size;
    ctx->fixed_size = info.fixed_size;
    ctx->component_count = info.component_count;
    ctx->bitmap_words = info.bitmap_words;
    ctx->format = info.format;

    bedrock_save(ctx->save_area);
    return 0;
}
\end{verbatim}
\endgroup\par

The component array is not part of the SAVE image. It is ordinary runtime metadata derived from CPUID so the
context owner can inspect, copy, or migrate extension components without hard-coding implementation-specific
offsets.
